\section{Exploratory Physical Analysis}\label{app:physics}
As a secondary question, we ask whether cross-body variability in control demand
differs between dynamic and quasi-static motions. The original MorphMimic checkpoint is
replayed over the 150 development clips and all 128 trained bodies (19,200
pairs, zero recorded failed episodes under this analysis's own diagnostic).
Keyword matching on the inherited captions assigns running, jumping, and
kicking to a dynamic group (25 clips) and standing and sitting to a
quasi-static group (37 clips); the other 88 clips are excluded, and the labels
have not been manually validated. After subtracting each clip's across-body
mean, Table~\ref{tab:physical_analysis} reports the standard deviation of the
pooled residuals, with 95\% intervals that resample clips as clusters (10,000
resamples). Torque, power, and jerk describe the simulated controller, not
calibrated human biomechanics.

\begin{table}[t]
\centering
\small
\caption{Cross-body control-demand variability, dynamic versus quasi-static clips.}\label{tab:physical_analysis}
\begin{tabularx}{\linewidth}{@{}Yrr@{}}
\toprule
Metric & Dynamic std [95\% CI] & Quasi-static std [95\% CI] \\
\midrule
Normalized torque (frac.\ effort limit) & $1.93$ [$1.71$, $2.21$] $\times10^{-4}$ & $2.06$ [$1.86$, $2.30$] $\times10^{-4}$ \\
Mass-normalized power & 0.292 [0.220, 0.370] & 0.259 [0.210, 0.323] \\
Normalized jerk & 350.2 [283.3, 416.9] & 403.2 [303.3, 513.5] \\
\bottomrule
\end{tabularx}
\end{table}

The spreads differ modestly and their separate intervals overlap, which is not
a test of the class-level difference. A direct dynamic-minus-quasi-static
contrast with clip-level resampling or permutation has not been estimated, so
the analysis supports descriptive comparisons only, not a positive or negative
interaction result.
